% ======================================================================
% SECTION 1: INTRODUCTION
% Location in paper: title page, continues onto page 2 above the table of
% contents. Three thematic blocks (a), (b), (c) - each 3 to 6 paragraphs.
% ======================================================================

[PROCESSING INSTRUCTION (a) - FDA April 28, 2026 News, 4 to 6 paragraphs:
Open the introduction with a focused account of the FDA's 28 April 2026
Real-Time Clinical Trials (RTCT) announcement. Read the full source file
at new-trial/national-24-7-trial/FDA-April-2026/FDA\_RealTime\_Clinical\_Trials.md
and quote, where appropriate, FDA Commissioner Marty Makary's exact framing
("For 60 years, we've been conducting clinical trials in the same way") and
Chief AI Officer Jeremy Walsh's quote on the transformative potential of
real-time trials. Identify the two announced proof-of-concept trials by
name: AstraZeneca's TRAVERSE Phase 2 in mantle cell lymphoma at MD Anderson
and the University of Pennsylvania, and Amgen's STREAM-SCLC Phase 1b in
limited-stage small cell lung carcinoma. Explain Paradigm Health's role as
the validated technical conduit for real-time signal sharing. Describe the
Request for Information for the upcoming pilot program. Close this block
by stating that the FDA's path to "continuous trials" - eliminating the
hiatus between phases - sets the regulatory stage on which the four
simulations in this paper operate. Be respectful of the FDA throughout;
frame the agency as enabling improvements rather than as the limiting
factor. Cite \cite{fda2026realtime}.]

[PROCESSING INSTRUCTION (b) - Current AI Patient Prediction Capabilities,
6 to 9 paragraphs: Read the following files in full and summarize the
current state of oncology AI patient prediction in measured, technical
prose:
- new-trial/national-24-7-trial/Background-A/README.md
- new-trial/national-24-7-trial/Background-A/chunk\_01\_baseline\_and\_short\_horizon.md
- new-trial/national-24-7-trial/Background-A/chunk\_02\_multimodal\_and\_limitations.md
- new-trial/national-24-7-trial/Background-A/chunk\_03\_bibtex.md
- new-trial/national-24-7-trial/Background-B/README.md
- new-trial/national-24-7-trial/Background-B/chunk\_01\_baseline\_and\_prediction\_domains.md
- new-trial/national-24-7-trial/Background-B/chunk\_02\_response\_metrics\_conclusions.md
- new-trial/national-24-7-trial/Background-B/chunk\_03\_bibtex\_references.md

Cover at least the following five domains, in order, with one short
paragraph per domain:

(b.1) Architectural reality: today's oncology prediction AI is a
heterogeneous collection of narrow supervised models (ridge Cox, SVM,
random survival forests, XGBoost, elastic net), not general-purpose
reasoning systems. Cite \cite{Smiley2025OncologyMLReportingQuality} and
\cite{Huang2025CancerSurvivalMetaAnalysis} for the null result that ML
provides SMD = 0.01 (95\% CI -0.01 to 0.03) over Cox regression on real-
world structured survival data.

(b.2) Short-horizon EHR-based mortality prediction: Manz 2020 with AUC
0.89 in 24,582 patients, PPV 45.2\% at 40\% threshold, and 4 to 8 day
prospective lead time. Patient-reported-outcome-only models (Sidey-Gibbons
2022) reach AUROC 0.69 to 0.76. Cite \cite{Manz2020Oncology180DayMortality}
and \cite{SideyGibbons2022OvarianPROMortality}.

(b.3) Adverse-event and hospitalization prediction: SHIELD-RT (Hong 2020)
as the gold-standard prospective RCT - 45\% reduction in acute-care events,
48\% reduction in costs - externally validated by Elia 2025 across 22,000
courses at AUROC 0.756 to 0.770. Add fluoropyrimidine cardiotoxicity (Li
2022 XGBoost AUC 0.816), neoadjuvant breast cancer toxicity (Cai 2024
AUROC reaching 0.75 with dose-intensity features), and the 2026 head and
neck pooled AUROC 0.76 (Ugwu 2026). Cite \cite{Hong2020SHIELDRT},
\cite{Elia2025SHIELDRTExternalValidation},
\cite{Li2022FluoropyrimidineCardiotoxicity},
\cite{Cai2024BreastCancerToxicityPrediction}, and
\cite{Ugwu2026HeadNeckRadiationToxicityMeta}.

(b.4) Medium and long-horizon survival prediction: AIM-LCpro (Li 2025)
multimodal NSCLC C-index 0.785 to 0.804 internal and 0.693 to 0.749
external; PROGPATH (Yuan 2025) pancancer ViT foundation model C-index
0.713 to 0.805 across 17 external cohorts; lung-cancer image-AI meta
sensitivity 0.83 / specificity 0.83 / AUC 0.90 (Yuan 2025 image meta).
Cite \cite{Li2025AIMLCpro}, \cite{Yuan2025PROGPATH}, and
\cite{Yuan2025LungImageMeta}.

(b.5) Response and progression prediction under immune checkpoint
inhibitors: SCORPIO routine-blood model median time-dependent AUC 0.763
hold-out and 0.725 external (Yoo 2025); Vanguri 2022 multimodal radiology
+ pathology + genomics AUC 0.80; Captier 2025 late-fusion C-index 0.75
OS, AUC 0.81 1-year death; PBMF contrastive biomarker discovery (Arango-
Argoty 2025); MMF in HCC (Xu 2025); Zhao 2023 anti-PD-1 radiomics in
breast cancer; Iivanainen 2021 ePRO + XGBoost for irAE detection. Cite
\cite{Yoo2025SCORPIO}, \cite{Vanguri2022MultimodalPDL1NSCLC},
\cite{Captier2025MultimodalNSCLCImmunotherapy}, \cite{ArangoArgoty2025PBMF},
\cite{Xu2025MMF}, \cite{Zhao2023BreastRadiomics}, and \cite{Iivanainen2021irAE}.

(b.6) Reporting and methodological floor: TRIPOD+AI (Collins 2024) and
CREMLS (El Emam 2024) define the reporting standards against which any new
predictive system - including the LLM-based simulations in this paper -
should be evaluated. Cite \cite{Collins2024TRIPODAI} and
\cite{ElEmam2024CREMLS}.

After the five technical paragraphs, add one short paragraph that names
the operational gaps left by the current baseline: retrospective bias,
heterogeneous endpoints, narrow per-task models, no real-time signal
sharing across sites and sponsors and FDA, and no contextual repository-
level reasoning across protocol, EHR, robotics, and regulatory documents
in a single inference pass.]

[PROCESSING INSTRUCTION (c) - Transition to the Four Simulations, 2
paragraphs: Transition from the gap analysis above into a one-paragraph
preview of the four author simulations and the central thesis: Claude Code
Opus 4.7 Max with a 1M token context window can hold protocol, regulatory,
robotic, EHR, and ASCII-diagram content in a single inference pass and
emit hour-by-hour patient predictions, safety signals, and forecasts at a
cadence and depth that no current narrow supervised model in oncology
trial practice has demonstrated. The second paragraph names the four
simulations exactly as listed in paper/README.md (Continuous RTCT;
Single-patient 10-stage journey; 24-hour autonomous sponsor; 168-hour
7-day extension with local verification on Core i5-6200U / 4 GB RAM) and
the Zenodo DOIs (10.5281/zenodo.19176370, 10.5281/zenodo.19119939, and
10.5281/zenodo.19396256 covering both sponsor simulations). Note that
the on-screen Table of Contents follows directly after this introduction
and on a fresh page. Cite \cite{kawchak_2026_19176370},
\cite{kawchak_2026_19119939}, and \cite{kawchak_2026_19396256}.]
